% GENERATED FILE. Edit canonical lessons, then run npm run generate:books.
% canonical-source-hash: fnv1a64:84d0122f1b84fe57
% canonical-lessons: HI-C359-mudda, HI-C359-samjhauta, HI-C359-vivad, HI-C359-ulajhna, HI-C359-cherna

\chapter{The Issue at Hand}
\label{ch:hi-the-issue-at-hand}
\hlchaptermodality{359}

\begin{chapteropening}
\textbf{By the end of this chapter:} \emph{I can say what a discussion is about and how it was settled.}

Complete the last lesson of chapter 359: I can say what a discussion is about and how it was settled.
\end{chapteropening}

\section[\texorpdfstring{\hi{मुद्दा} (muddā)}{मुद्दा (muddā)}]{\hi{मुद्दा} (muddā) — an issue, a point at issue}

\label{lesson:HI-C359-mudda}



\begin{quote}
Before the new one: say the Hindi for greed, then the Hindi for to make known.
\end{quote}

\subsection*{You'll want to know: \hi{मुद्दा}}
\textbf{\hi{मुद्दा}} — \emph{muddā} — \textquotedblleft{}an issue, a point at issue\textquotedblright{}.

\begin{quote}
\hi{आज} \hi{की} \hi{बैठक} \hi{में} \hi{पानी} \hi{सबसे} \hi{बड़ा} \hi{मुद्दा} \hi{था।}
\end{quote}

\begin{quote}
\emph{āj kī baiṭhak meṁ pānī sabse baṛā muddā thā.}
\end{quote}

\begin{quote}
Water was the biggest issue at today's meeting.
\end{quote}

Masculine.

\begin{cousinweb}
\textbf{\hi{मुद्दा}} (\emph{muddā}) is usually traced to Arabic \emph{muddaʿā}, \textquotedblleft{}that which is claimed\textquotedblright{}, which reached Hindi through Persian.
\end{cousinweb}

\begin{grammarlens}[title={What you've built}]
One more word for giving reasons and opinions.
\end{grammarlens}

\subsection*{Your turn}
\begin{itemize}
\raggedright
  \item \emph{Say it:} \emph{muddā}
  \item \emph{Say it:} \emph{muddā}, once more
  \item \emph{Say it:} \emph{āj kī baiṭhak meṁ pānī sabse baṛā muddā thā.} — Water was the biggest issue at today's meeting.
  \item \emph{Recall:} say the Hindi for greed, then the Hindi for to make known, then say \emph{muddā} again
\end{itemize}

\begin{tcolorbox}[breakable,colback=teal!4,colframe=teal!35!black,title={Before you move on}]
What does \hi{मुद्दा} mean? (\textquotedblleft{}An issue, a point at issue\textquotedblright{}.) Say it once more.
\end{tcolorbox}
\section[\texorpdfstring{\hi{समझौता} (samjhautā)}{समझौता (samjhautā)}]{\hi{समझौता} (samjhautā) — a compromise, a settlement}

\label{lesson:HI-C359-samjhauta}



\begin{quote}
Before the new one: say the Hindi for to make known, then the Hindi for an issue.
\end{quote}

\subsection*{You'll want to know: \hi{समझौता}}
\textbf{\hi{समझौता}} — \emph{samjhautā} — \textquotedblleft{}a compromise, a settlement\textquotedblright{}.

\begin{quote}
\hi{आख़िर} \hi{दोनों} \hi{परिवारों} \hi{में} \hi{समझौता} \hi{हुआ।}
\end{quote}

\begin{quote}
\emph{ākhir donõ parivāroṁ meṁ samjhautā huā.}
\end{quote}

\begin{quote}
In the end the two families reached a settlement.
\end{quote}

Masculine.

\begin{cousinweb}
\textbf{\hi{समझौता}} (\emph{samjhautā}) is a noun made on \hi{समझना}, \textquotedblleft{}to understand\textquotedblright{}, which CDIAL traces to Sanskrit \emph{saṃbudhyate}, \textquotedblleft{}understands\textquotedblright{} (CDIAL 12959): an understanding reached between two sides.
\end{cousinweb}

\begin{grammarlens}[title={What you've built}]
One more word for giving reasons and opinions.
\end{grammarlens}

\subsection*{Your turn}
\begin{itemize}
\raggedright
  \item \emph{Say it:} \emph{samjhautā}
  \item \emph{Say it:} \emph{samjhautā}, once more
  \item \emph{Say it:} \emph{ākhir donõ parivāroṁ meṁ samjhautā huā.} — In the end the two families reached a settlement.
  \item \emph{Recall:} say the Hindi for to make known, then the Hindi for an issue, then say \emph{samjhautā} again
\end{itemize}

\begin{tcolorbox}[breakable,colback=teal!4,colframe=teal!35!black,title={Before you move on}]
What does \hi{समझौता} mean? (\textquotedblleft{}A compromise, a settlement\textquotedblright{}.) Say it once more.
\end{tcolorbox}
\section[\texorpdfstring{\hi{विवाद} (vivād)}{विवाद (vivād)}]{\hi{विवाद} (vivād) — a dispute}

\label{lesson:HI-C359-vivad}



\begin{quote}
Before the new one: say the Hindi for an issue, then the Hindi for a compromise.
\end{quote}

\subsection*{You'll want to know: \hi{विवाद}}
\textbf{\hi{विवाद}} — \emph{vivād} — \textquotedblleft{}a dispute\textquotedblright{}.

\begin{quote}
\hi{दो} \hi{परिवारों} \hi{में} \hi{पानी} \hi{पर} \hi{विवाद} \hi{है।}
\end{quote}

\begin{quote}
\emph{do parivāroṁ meṁ pānī par vivād hai.}
\end{quote}

\begin{quote}
Two families have a dispute over water.
\end{quote}

Masculine.

\begin{cousinweb}
\textbf{\hi{विवाद}} (\emph{vivād}) is a tatsama, borrowed from Sanskrit \emph{vivāda}, \textquotedblleft{}a dispute, a contest\textquotedblright{} (Monier-Williams).
\end{cousinweb}

\begin{grammarlens}[title={What you've built}]
One more word for giving reasons and opinions.
\end{grammarlens}

\subsection*{Your turn}
\begin{itemize}
\raggedright
  \item \emph{Say it:} \emph{vivād}
  \item \emph{Say it:} \emph{vivād}, once more
  \item \emph{Say it:} \emph{do parivāroṁ meṁ pānī par vivād hai.} — Two families have a dispute over water.
  \item \emph{Recall:} say the Hindi for an issue, then the Hindi for a compromise, then say \emph{vivād} again
\end{itemize}

\begin{tcolorbox}[breakable,colback=teal!4,colframe=teal!35!black,title={Before you move on}]
What does \hi{विवाद} mean? (\textquotedblleft{}A dispute\textquotedblright{}.) Say it once more.
\end{tcolorbox}
\section[\texorpdfstring{\hi{उलझना} (ulajhnā)}{उलझना (ulajhnā)}]{\hi{उलझना} (ulajhnā) — to get tangled, to get bogged down}

\label{lesson:HI-C359-ulajhna}



\begin{quote}
Before the new one: say the Hindi for a compromise, then the Hindi for a dispute.
\end{quote}

\subsection*{You'll want to know: \hi{उलझना}}
\textbf{\hi{उलझना}} — \emph{ulajhnā} — \textquotedblleft{}to get tangled, to get bogged down\textquotedblright{}.

\begin{quote}
\hi{यह} \hi{धागा} \hi{हमेशा} \hi{उलझता} \hi{है।}
\end{quote}

\begin{quote}
\emph{yah dhāgā hameśā ulajhtā hai.}
\end{quote}

\begin{quote}
This thread always gets tangled.
\end{quote}

Past (it …): \hi{वह} \hi{धागा} \hi{उलझ} \hi{गया}, \emph{vah dhāgā ulajh gayā}.

\begin{cousinweb}
\textbf{\hi{उलझना}} (\emph{ulajhnā}) is a tadbhava, inherited from Sanskrit \emph{uparundhati}, \textquotedblleft{}obstructs\textquotedblright{} (CDIAL 2221).
\end{cousinweb}

\begin{grammarlens}[title={What you've built}]
One more word for giving reasons and opinions.
\end{grammarlens}

\subsection*{Your turn}
\begin{itemize}
\raggedright
  \item \emph{Say it:} \emph{ulajhnā}
  \item \emph{Say it:} \emph{ulajhnā}, once more
  \item \emph{Say it:} \emph{yah dhāgā hameśā ulajhtā hai.} — This thread always gets tangled.
  \item \emph{Recall:} say the Hindi for a compromise, then the Hindi for a dispute, then say \emph{ulajhnā} again
\end{itemize}

\begin{tcolorbox}[breakable,colback=teal!4,colframe=teal!35!black,title={Before you move on}]
What does \hi{उलझना} mean? (\textquotedblleft{}To get tangled, to get bogged down\textquotedblright{}.) Say it once more.
\end{tcolorbox}
\section[\texorpdfstring{\hi{छेड़ना} (cheṛnā)}{छेड़ना (cheṛnā)}]{\hi{छेड़ना} (cheṛnā) — to raise (a topic); to tease}

\label{lesson:HI-C359-cherna}



\begin{quote}
Before the new one: say the Hindi for a dispute, then the Hindi for to get tangled.
\end{quote}

\subsection*{You'll want to know: \hi{छेड़ना}}
\textbf{\hi{छेड़ना}} — \emph{cheṛnā} — \textquotedblleft{}to raise (a topic); to tease\textquotedblright{}.

\begin{quote}
\hi{बैठक} \hi{में} \hi{अरुण} \hi{ने} \hi{पानी} \hi{का} \hi{मुद्दा} \hi{छेड़ा।}
\end{quote}

\begin{quote}
\emph{baiṭhak meṁ Aruṇ ne pānī kā muddā cheṛā.}
\end{quote}

\begin{quote}
At the meeting Arun brought up the water issue.
\end{quote}

Past, with \hi{ने} (he or she …): \hi{उसने} \hi{वही} \hi{बात} \hi{फिर} \hi{छेड़ी}, \emph{usne vahī bāt phir cheṛī}.

\begin{cousinweb}
\textbf{\hi{छेड़ना}} (\emph{cheṛnā}) is inherited from \emph{*chēḍ-}, a form CDIAL reconstructs from its descendants and files under \emph{*ciḍ}, \textquotedblleft{}be angry\textquotedblright{} (CDIAL 4794). CDIAL glosses the Hindi verb \textquotedblleft{}to touch, to irritate\textquotedblright{}; raising a topic is a sense it does not give.
\end{cousinweb}

\begin{grammarlens}[title={What you've built}]
One more word for giving reasons and opinions.
\end{grammarlens}

\subsection*{Your turn}
\begin{itemize}
\raggedright
  \item \emph{Say it:} \emph{cheṛnā}
  \item \emph{Say it:} \emph{cheṛnā}, once more
  \item \emph{Say it:} \emph{baiṭhak meṁ Aruṇ ne pānī kā muddā cheṛā.} — At the meeting Arun brought up the water issue.
  \item \emph{Recall:} say the Hindi for a dispute, then the Hindi for to get tangled, then say \emph{cheṛnā} again
\end{itemize}

\begin{tcolorbox}[breakable,colback=teal!4,colframe=teal!35!black,title={Before you move on}]
What does \hi{छेड़ना} mean? (\textquotedblleft{}To raise (a topic); to tease\textquotedblright{}.) Say it once more.
\end{tcolorbox}
